\documentclass[12pt,letterpaper]{article}
% Common article worksheet layout. Load after \documentclass.
\usepackage[margin=0.5in,top=0.85in,headheight=0pt,headsep=0pt,footskip=0pt]{geometry}
\usepackage{amsmath,amssymb}
\usepackage{tikz,graphicx}
\definecolor{mtuGreen}{RGB}{0,79,42}
\definecolor{mtuGreenLight}{RGB}{60,133,90}
\definecolor{mtuGreenPale}{RGB}{223,236,227}
\newcommand{\vect}[1]{\mathbf{#1}}
\setlength{\parindent}{0pt}
\setlength{\parskip}{8pt}
\pagestyle{empty}
% Keep all marks inside a half-inch print margin, including their strokes.
% Odd pages are left-aligned; even pages are right-aligned.
\AddToHook{shipout/foreground}{%
  \begin{tikzpicture}[remember picture,overlay]
    \ifodd\value{page}
      \node[anchor=north west,draw=black,rectangle,line width=1.5pt,inner xsep=3pt,inner ysep=1pt,font=\fontsize{16}{16}\selectfont\bfseries]
        at ([xshift=0.52in,yshift=-0.52in]current page.north west) {\thepage};
    \else
      \node[anchor=north east,draw=black,rectangle,line width=1.5pt,inner xsep=3pt,inner ysep=1pt,font=\fontsize{16}{16}\selectfont\bfseries]
        at ([xshift=-0.52in,yshift=-0.52in]current page.north east) {\thepage};
    \fi
    % First-page marker matches the dimensions of the boxed page number 1.
    \ifnum\value{page}=1\relax
      \node[anchor=north east,draw=black,fill=black,rectangle,line width=1.5pt,inner xsep=3pt,inner ysep=1pt,font=\fontsize{16}{16}\selectfont\bfseries]
        at ([xshift=-0.52in,yshift=-0.52in]current page.north east) {\phantom{1}};
    \fi
  \end{tikzpicture}%
}
% Fixed name-box dimensions and page position shared by every worksheet.
\newcommand{\WorksheetNameBox}{%
  \begingroup\setlength{\fboxsep}{3pt}\setlength{\fboxrule}{0.4pt}%
  \fbox{\rule[-0.30in]{0pt}{0.42in}\makebox[3.3in][l]{Name:}}%
  \endgroup}
\newcommand{\WorksheetHeading}[1]{%
  \noindent{\small\textbf{MA3521 --- #1}}\hfill\WorksheetNameBox\par}
% A blank entry with fixed writing room, usable inside a matrix.
\newcommand{\AnswerEntry}{\rule[-0.17in]{0pt}{0.47in}\makebox[0.67in]{}}
\newcommand{\AnswerVector}{\begin{pmatrix}\AnswerEntry\\\AnswerEntry\\\AnswerEntry\end{pmatrix}}
\newcommand{\AnswerMatrix}{\begin{pmatrix}
  \AnswerEntry&\AnswerEntry&\AnswerEntry\\
  \AnswerEntry&\AnswerEntry&\AnswerEntry\\
  \AnswerEntry&\AnswerEntry&\AnswerEntry
\end{pmatrix}}

\usepackage{bm}
% The worksheet header supplies \vect. The MTU theme supplies this \mat definition.
\newcommand{\mat}[1]{\bm{#1}}
\newcommand{\SystemBlank}[2]{%
  \vcenter{\hbox{$\boxed{\rule{0pt}{#2}\makebox[#1]{}}$}}}
\newcommand{\VectorBlank}{\SystemBlank{3.2in}{0.80in}}
\newcommand{\MatrixBlank}{\SystemBlank{1.8in}{0.80in}}
\newcommand{\TrigMatrixBlank}{\SystemBlank{3.0in}{1.60in}}
\begin{document}
\noindent\makebox[0pt][l]{\raisebox{-\baselineskip}[0pt][0pt]{\small Do not row reduce or solve.}}%
\WorksheetHeading{In-Class Worksheet Monday/Tuesday}
\begin{enumerate}
\item
\begingroup\setlength{\arraycolsep}{3pt}
$\vect{y}'=\mat{A}\vect{y}+e^{-2t}\begin{pmatrix}-3\\-5\\-1\end{pmatrix}\quad \mat{A}=\begin{pmatrix}1&-2&0\\2&1&0\\0&0&-1\end{pmatrix}\quad \vect{y}(0)=\begin{pmatrix}2\\3\\1\end{pmatrix}$
\par Eigenpairs $1+2i,\,\begin{pmatrix}1\\-i\\0\end{pmatrix}\quad 1-2i,\,\begin{pmatrix}1\\i\\0\end{pmatrix}\quad -1,\,\begin{pmatrix}0\\0\\1\end{pmatrix}$
\endgroup
\begin{enumerate}
\item 
$\vect{y}_c(t)=\VectorBlank$
\item A suitable UC form is
\newline
$\vect{y}_p(t)=\VectorBlank$
\item To determine the UC coefficients 
you would row reduce \newline
$\mat{M}_{\text{\textsc{uc}}}=\MatrixBlank$
\item To satisfy the IC you would row reduce
\par\noindent $\mat{M}_{\text{\textsc{ic}}}=\MatrixBlank$\par
\end{enumerate}
\vfill
Not on exam  $\vect{y}(t)=\VectorBlank$

\newpage
\item
\begingroup\setlength{\arraycolsep}{3pt}
$\vect{y}'=\mat{A}\vect{y}+\cos(2t)\begin{pmatrix}-3\\-5\\-1\end{pmatrix}\quad \mat{A}=\begin{pmatrix}-1&-2&1\\-2&-3&3\\-1&-3&2\end{pmatrix}\quad \vect{y}(0)=\begin{pmatrix}2\\3\\1\end{pmatrix}$
\par Eigenpairs $i,\,\begin{pmatrix}1\\-i\\1-i\end{pmatrix}\quad -i,\,\begin{pmatrix}1\\i\\1+i\end{pmatrix}\quad -2,\,\begin{pmatrix}1\\1\\1\end{pmatrix}$
\endgroup
\begin{enumerate}
\item 
$\vect{y}_c(t)=\VectorBlank$
\item A suitable UC form is
\newline
$\vect{y}_p(t)=\VectorBlank$
\item To determine the UC coefficients 
you would row reduce \newline
$\mat{M}_{\text{\textsc{uc}}}=\SystemBlank{2.7in}{1.20in}$
\item To satisfy the IC you would row reduce
\par\noindent $\mat{M}_{\text{\textsc{ic}}}=\MatrixBlank$\par
\end{enumerate}
\vfill
Not on exam  $\vect{y}(t)=\VectorBlank$

\end{enumerate}
% Adapted from Yprime=AYplusFPartTwo.tex. No row reduction required.
% Front: characteristic polynomial (lambda+1)((lambda-1)^2+4).
% Front UC form a exp(-2t). With columns a_1,a_2,a_3,
% UC augmented matrix [-3 2 0 | -3; -2 -3 0 | -5; 0 0 -1 | -1].
% Real basis exp(t)(cos(2t),sin(2t),0),
% exp(t)(-sin(2t),cos(2t),0), exp(-t)(0,0,1).
% IC matrix [I | (2,3,1)^T-a] for that basis.
% Back: characteristic polynomial (lambda+2)(lambda^2+1).
% UC form a cos(2t)+b sin(2t). Columns a_1,a_2,a_3,b_1,b_2,b_3.
% Cosine equations first, then sine equations. The augmented matrix is
% [1 2 -1 2 0 0 | -3; 2 3 -3 0 2 0 | -5; 1 3 -2 0 0 2 | -1;
%  -2 0 0 1 2 -1 | 0; 0 -2 0 2 3 -3 | 0; 0 0 -2 1 3 -2 | 0].
% One real complementary basis is Re(exp(it)v), -Im(exp(it)v),
% exp(-2t)(1,1,1), where v=(1,-i,1-i).
% For this basis, IC matrix [1 0 1 | 2-a_1; 0 1 1 | 3-a_2; 1 1 1 | 1-a_3].
\end{document}
